\NeedsTeXFormat{LaTeX2e}
\documentclass[a4paper,11pt,
headsepline,           % Linie zw. Kopfzeile und Text
oneside,               % einseitig
pointlessnumbers,      % keine Punkte nach den letzten Ziffern in Überschriften
bibtotoc,              % LV im IV
DIV=15,               % Satzspiegel auf 15er Raster, schmalere Ränder   
BCOR15mm               % Bindekorrektur
%,draft
]{scrartcl}
\KOMAoptions{DIV=last} % Neuberechnung Satzspiegel nach Laden von Paket helvet

\pagestyle{headings}
\usepackage{blindtext}

\usepackage[utf8]{inputenc}
\usepackage[T1]{fontenc}

\usepackage[scaled]{helvet}
\renewcommand{\familydefault}{\sfdefault} 

\usepackage{graphicx}

\usepackage{csquotes}
\usepackage[backend=bibtex8]{biblatex}
\bibliography{bibliography}

\usepackage{tabularx} 

\usepackage{url}              % \url{http://...} in Schreibmaschinenschrift
\usepackage{color}            % zum Setzen farbigen Textes

\usepackage{amssymb, amsmath} % Pakete für Mathe-Umgebungen und -Symbole

\usepackage{setspace}         % Paket für div. Abstände, z.B. ZA
\setlength{\parindent}{0pt}   % kein linker Einzug der ersten Absatzzeile
\setlength{\parskip}{1.4ex plus 0.35ex minus 0.3ex} % Absatzabstand, leicht variabel

\setcounter{tocdepth}{3}      % ist Standard

\usepackage{listings}
\lstset{language=C++,
        basicstyle=\ttfamily,
        keywordstyle=\color{blue}\ttfamily,
        stringstyle=\color{red}\ttfamily,
        commentstyle=\color{magenta}\ttfamily,
        morecomment=[l][\color{green}]{\#},
        breaklines=true,
        numbers=left,
        stepnumber=1,
        morekeywords={constexpr, concept}
}


\usepackage{todonotes}

\newcommand{\fullname}{Paul Nykiel}
\newcommand{\email}{paul.nykiel@uni-ulm.de}
\newcommand{\titel}{Advanced Concepts of Metaprogramming for the Implementation of a CPU Emulator Using Templates}
\newcommand{\jahr}{2020}
\newcommand{\matnr}{941496}
\newcommand{\gutachter}{Dr.\,Borchert}

% hier die Fakultät auswählen
\newcommand{\fakultaet}{Mathematik und Wirtschaftswissenschaften}

% hier das Institut einsetzen
\newcommand{\institut}{Institut für Numerische Mathematik}

% Informationen, die LaTeX in die PDF-Datei schreibt
\pdfinfo{
  /Author (\fullname)
  /Title (\titel)
  /Producer     (pdfeTex 3.14159-1.30.6-2.2)
  /Keywords ()
}

\usepackage{hyperref}
\hypersetup{
pdftitle=\titel,
pdfauthor=\fullname,
pdfsubject={Diplomarbeit},
pdfproducer={pdfeTex 3.14159-1.30.6-2.2},
colorlinks=false,
pdfborder=0 0 0	% keine Box um die Links!
}

% Trennungsregeln
\hyphenation{Sil-ben-trenn-ung}

\begin{document}

% Titelseite
\thispagestyle{empty}
\begin{addmargin*}[4mm]{-10mm}

\includegraphics[height=1.8cm]{images/unilogo_bild}
\hfill
\includegraphics[height=1.8cm]{images/unilogo_wort}\\[1em]

{\footnotesize
%{\bfseries Universität Ulm} \textbar ~89069 Ulm \textbar ~Germany
\hspace*{115mm}\parbox[t]{35mm}{\bfseries Fakultät für\\
\fakultaet\\
% TODO hier Institut anpassen
\mdseries \institut}\\[2cm]

\parbox{140mm}{\bfseries \LARGE \titel}\\[2.5em]
{\footnotesize Hausarbeit zur Vorlesung "Objektorientierte Programmierung mit C++
"}\\[3em]

{\footnotesize \bfseries Vorgelegt von:}\\
{\footnotesize \fullname\\ \email}\\ \matnr\\[2em]
{\footnotesize \bfseries Gutachter:}\\                     
{\footnotesize \gutachter}\\[2em]
{\footnotesize \jahr}
}
\end{addmargin*}


% Impressum
\clearpage
\thispagestyle{empty}
{ \small
  \flushleft
  Version \today \\\vfill
  \copyright~\jahr~\fullname\\[0.5em]

This report is free software: you can redistribute it and/or modify
it under the terms of the GNU General Public License as published by
the Free Software Foundation, either version 3 of the License, or
(at your option) any later version.

This report is distributed in the hope that it will be useful,
but WITHOUT ANY WARRANTY; without even the implied warranty of
MERCHANTABILITY or FITNESS FOR A PARTICULAR PURPOSE.  See the
GNU General Public License for more details.

You should have received a copy of the GNU General Public License
along with this program.  If not, see \href{https://www.gnu.org/licenses}{https://www.gnu.org/licenses}.\\

  Satz: PDF-\LaTeXe
}

\newpage

% ab hier Zeilenabstand etwas größer 
\setstretch{1.2}

\tableofcontents

\newpage

\section{Problem Statement}
%1. Was ist das Problem / die Aufgabe

%Hierbei würde
%ich auf ein Projekt aufbauen, an dem ich schon die letzten Wochen
%arbeite. Es geht um die Implementierung eines einfachen
%CPU-Emulators (die Architektur ist teilweise von MIPS inspiriert)
%rein über das C++-Typsystem (zur verbesserten Typsicherheit werden
%außerdem Concepts genutzt). Programme sind jeweils ein eigener Typ
%und werden vollständig zur Compile-Time ausgeführt, dadurch kann gut
%gezeigt werden dass bereits das C++-Typsystem Turing-Vollständig
%ist. Den aktuellen Stand der Entwicklung finden sie unter:
%https://github.com/aul12/TemplateCpu. Im Rahmen der Arbeit würde ich
%zuerst die genutzten Konzepte vorstellen, und dann erklären wie ein
%Programm ausgeführt wird, für das Softwarepaket würde ich das
%Programm um eine ordentliche Dokumentation erweitern und eine Turing
%Maschine als Program für den CPU-Emulator implementieren.

In comparison to most other widely adopted languages the C++ language provides a very expressive type system \cite{concepts05}. Custom
types (\lstinline{structs} and \lstinline{classes}) do not only consist of other types but can also depend on other
types and values. This yields types which can fundamentaly change their behaviour based on other types \cite[Chapter~13.3]{std}. As a result
it is possible to implement algorithms which only consist of the type system and dependent types, this concept is known
as metaprogramming. In the case of C++ this metaprogramming is turing complete, this means that all algorithms than
can be computed can also be computed by the C++ type system \cite{TuringComputability}.

In this report the implementation of a CPU emulator using only the C++ type system is presented. This emulator
accepts programs, which are custom types that behave like instructions, and executes these during the compilation of
the program. These programs can be, in theory, arbitrary and thus all computable algorithms can be implemented. As a
demonstration of the Turing completeness a Turing machine will be implemented, the Turing machine is a concept which
is used for the definition of Turing computability \cite{Turing1936}.

This report is structured into three parts: first the necessary concepts and data types for implementing more complex 
programs using metaprogramming are introduced. In the second part the implementation of the CPU emulator is described
and an overview on how to implement programs for the emulator is given. In the last part a conclusion is formulated and
possible improvements in the future are listed.


\section{Concepts}
%2. Wie bin ich generell da ran gegangen? Was für Konzepte habe ich mir überlegt, welche Strukturen, patterns...
\subsection{Template Specialization}
Metaprogramming in C++ is primarily based on the concept of template-specializiation: for a specific template instantiation
a conrete implementation of the class or function can be given. This makes it possible to implement control structures
using the type system. 

Listing \ref{lst:template_specialization} implements a logical not using structs-template with specialization.
\todo{Disclaimer ! geht auch}

\lstinputlisting[caption={Template specialization}, label=lst:template_specialization]{template_specialization.cpp}

For struct- and class-templates, but not for functions-templates, it is also possible to do partial template 
specialization, if a struct of class depends on more than one type or value a specialization is allowed to define
only a subset of the arguments and leave other arguments as template arguments. Listing
\ref{lst:partial_template_specialization} uses partial template specialization to define a struct template which
behaves similar to \lstinline{std::enable_if}: if the first argument is true the class defines a type, if the
argument is false this type is not defined.

\lstinputlisting[caption={Partial template specialization}, label=lst:partial_template_specialization]
    {partial_template_specialization.cpp}

\subsection{Compile-Time Datatypes}
For implementing algorithms it is not only necessary to be able to implement the program logic and control flow
it is equally important to define data types to work with. When implementing a CPU emulator the required data type
is primarily a sequential container for representing registers, memory and the program itself. 
As the STL does not provide any data types suitable for compile time manipulation a custom container had to be implemented.

The implementation of the sequential container is based on singly linked list. 
This not the obvious choice when implementing sequential containers but array manipulation is not possible during 
compile time, thus a different implementation of a sequential container had to be chosen.

The definition of the list is similar to a typical singly linked list definition in normal run-time C++. Every
element consists of a value and the next element, for the last element the next element is of type \lstinline{ListEnd}.
Listing \ref{lst:typelist} is the definition of the type list, that is a list in which every element is a type, not
a value. This type of list is used for the program of the CPU emulator, as the instructions are types not values.

\lstinputlisting[caption={Type List}, label={lst:typelist}]{type_list.cpp}

Building upon the type list we can define a value list as a type list of \lstinline{ValueContainer}s. Listing
\ref{lst:valuecontainer} gives the definition of the value container struct. This struct contains a value but
is a type itself, thus a type list can be used. This scheme makes it possible to reuse most of the type list
implementation for the value list.

\lstinputlisting[caption={Value Container}, label={lst:valuecontainer}]{value_container.cpp}


\todo[inline]{Helper Functions}

\subsection{Concepts}
\todo[inline]{Motivation}

\todo[inline]{Usage}

\subsection{Limitations}
\todo[inline]{Max Depth}


\section{Implementation}
%3. Dann natürlich die Lösung vorstellen, gerne auch interessante Code-Schnipsel zeigen

\subsection{Execution of Code}
For the emulation the most important data structure is the CPU state. The state consists of the program, the registers,
the memory and the program counter among others. During each CPU state the next state is determined based on the 
current state and the current instruction.
The emulator terminates once the program counter is equal to the length of the program,
the emulator then returns a result containing the program counter, the current content
of the registers and the content of the memory.

The program, the registers and the memory are all defined using  compile time lists.
For the program the list is a type list as the instructions are represented by different
types. The registers are a value list, the value type can be configured (see the
\texttt{config.hpp} header), and is \lstinline{int} by default. The registers are not
accessed directly using indizes, but via an enum class which provides
more readable names for the registers. The memory is a value list with the same type
as the registers to avoid problems related to conversions. The size of the memory
can be configured, depending on the requirements but the limitations of the compiler 
(see Section \ref{sec:limitations}), need to be respected.

For the execution of the instructions a second struct called \lstinline{ExecuteInstr} is defined. 
This struct fetches the instruction based on the program counter and the program. Then
the instruction is executed and the program counter, the registers and the memory is
updated based on the result of the execution.

The logic used for the execution of instructions is show in Listing \ref{lst:cpu}.

\lstinputlisting[caption={Implementation of the main emulator code}, label={lst:cpu}]{code/cpu.cpp}

\subsection{Instructions}

For programming the CPU emulator a total of 28 instructions have been implemented. They consist of arithmetic and logic
operations, load and store operations and jumps, both conditional and unconditional. The instructions are primarily 
influenced by the MIPS architecture \todo{cite}, similar to MIPS the instructions are based
on a Load/Store architecture \todo{cite}. For the full list of instructions refer to the \texttt{README.md} file provided with the code. 

For the execution there is a third struct called \lstinline{InstrImpl} which is
specialized for all instructions, depending on the instruction the program counter,
registers and memory is then manipulated in each template specizalization.

Listing \ref{lst:add} shows the implementation of the \lstinline{Add} instruction.
For this the \lstinline{InstrImpl} struct is specialized for the respective type.
This struct then defines three members. The first member are the registers after the
instruction is executed. For the execution first the two operands are fetched using
\lstinline{GetVal}, then the values are summed up and the result is stored into a
register using \lstinline{SetVal}. The member containing the next programm counter
is simply calculated by incrementing the program counter by one. The memory is not
changed by the instruction, thus it is the same as before.

\lstinputlisting[caption={Implementation of the \lstinline{Add} instruction},label={lst:add}]{code/add.cpp}

Listing \ref{lst:jmp} shows the implementation of the \lstinline{Jump} instruction.
Similar to the \lstinline{Add} instruction this is achieved by specializing the 
\lstinline{InstrImpl} struct. In the implementation both the registers and the memory
are not changed. Only the program counter is changed, for this the address is fetched
from the respective register using \lstinline{GetVal} and then the program counter is
set to this value.

\lstinputlisting[caption={Implementation of the \lstinline{Jump} instruction},label={lst:jmp}]{code/jmp.cpp}

\todo[inline]{Load/Store?}

\subsection{Definining and Running a Program}


\todo[inline]{Printing results}

\subsection{Debugger}
\todo[inline]{Debugger}


\section{Conclusion}
%4. Ausblick, was könnte man erweitern (warum ist das idealerweise einfach? => Weil dein code halt so krass gut aufgebaut ist...), oder auch anders machen?

\todo[inline]{Fazit: Applications}

This report explained how a CPU emulator can be implemented using metaprogramming techniques.
The necessary basics, such as compile time data types and concepts, have been explained 
in the first section. The second section explained the structure of the emulator and how
programs are executed by the emulator. Additionally a program for emulating a turing
machine has been implemented for the CPU emulator to prove that the emulator, and thus
C++ metaprogramming is turing complete. It is important to consider that while in theory
arbitrary programs can be executed in reality the programs are limited by the compiler.

\subsection{Future Improvements}
If necessary more instructions can be added easily in the future. A new instructions
requires the type declaration and the specialization of the \lstinline{InstrImpl}
struct with the implementation of the instruction. The current instruction set
is sufficient to implement all programs, but new instructions might reduce the
complexity of programs significantly.

Another improvement which would improve the usability of the emulator is the implementation
of input and output operations. Currently there is no easy way of doing this, there are
similar programs which implement I/O operations by recompiling the program for
every input \todo{src: https://jguegant.github.io/blogs/tech/meta-crush-saga.html,
https://github.com/Drahflow/compile-time-adventures} which could be implemented for
emulator but would require additional runtime code for recompiling the program.

Another usability improvement for the emulator, especially when writing large programs,
would be the addition of labels, like they are used in most assembly dialects.
Such as labels will be implemented in the future as part of an assembler, but are
currently not supported as they are not a feature provided by a CPU but by an assembler.


\newpage
\appendix

\nocite{Knappen2009}
\nocite{Mittelbach2005}
\nocite{Schlosser2014}
\nocite{Sturm2012}
\nocite{Voss2010}

\printbibliography

%\clearpage
%\thispagestyle{empty}

%Name: \fullname \hfill Matrikelnummer: \matnr \vspace{2cm}

%\minisec{Erklärung}

%Ich erkläre, dass ich die Arbeit selbständig verfasst und keine anderen als die angegebenen Quellen und Hilfsmittel verwendet habe.\vspace{2cm}

%Ulm, den \dotfill

%\hspace{10cm} {\footnotesize \fullname}

\end{document}
